% 分野別類題: 微分計算（逆三角関数）
% 逆三角関数の微分公式の導出（基本3問）と、合成・積を含む計算（やや複雑2問）。
% コンパイル: TEXINPUTS="../../../00_common:$TEXINPUTS" latexmk -xelatex problems.tex
\documentclass[a4paper,11pt]{article}
\usepackage{preamble}

\begin{document}

\kamokumei{類題演習 --- 微分計算（逆三角関数）}

\shijibun{以下の問いに答えよ。逆三角関数の微分公式を用いる場合は、その導出過程（逆関数の微分法）も示すこと。ただし $\arcsin x,\ \arccos x$ の値域はそれぞれ $\left[-\dfrac{\pi}{2}, \dfrac{\pi}{2}\right]$、$[0, \pi]$、$\arctan x$ の値域は $\left(-\dfrac{\pi}{2}, \dfrac{\pi}{2}\right)$ とする。}

\shomon{問1.}{$y = \arcsin x$（$-1 < x < 1$）について、逆関数の微分法を用いて $\dfrac{dy}{dx}$ を導出せよ。}

\shomon{問2.}{$y = \arctan x$ について、逆関数の微分法を用いて $\dfrac{dy}{dx}$ を導出せよ。}

\shomon{問3.}{次の関数を微分せよ。}
\[
y = \arccos(3x)
\]

\shomon{問4.}{次の関数を微分せよ。}
\[
y = \arctan\!\left(\frac{2x}{1-x^{2}}\right)
\]

\shomon{問5.}{次の関数を微分せよ。またその結果をできるだけ簡単な形にせよ。}
\[
y = x\arcsin x + \sqrt{1-x^{2}}
\]

\end{document}
